\documentclass{article}

% Margins and spacing
\usepackage[a4paper, margin=.7in]{geometry}
% Math packages
\usepackage{amsmath}
\usepackage{amssymb}
% Figures and graphics
\usepackage{graphicx}
\usepackage{caption}
% Placement of figures, tables, algorithms, etc.
\usepackage{float}
% Hyperlinks
\usepackage{hyperref}
% Reference section names
\usepackage{nameref}
% Algorithms
\usepackage{algorithm}
\usepackage{algpseudocode}
% Sus
\usepackage{listings}
\usepackage{xcolor}

%% Tips
% Fix spacing between lines
% \vspace{-\baselineskip}

% No section numbering
\setcounter{secnumdepth}{0}
% No page numbering
\pagenumbering{gobble}
% No paragraph indentation
\setlength{\parindent}{0pt}
% Removes caption prefixes (e.g., "Figure 1:") from figures
% \captionsetup[figure]{labelformat=empty}

% Define custom commands
%% Underlined hyperlinks
\newcommand{\uhref}[2]{\href{#1}{\underline{#2}}}
%% Bold-italic math symbols
\newcommand{\mathbi}[1]{\boldsymbol{#1}}
%% Caret symbol
\newcommand{\caret}{\text{\^{}}}
%% Operators
\DeclareMathOperator*{\argmin}{arg\,min}
\DeclareMathOperator*{\argmax}{arg\,max}
\DeclareMathOperator{\sgn}{sgn}
%% Manually indented states in algorithmic environment
\newcommand{\IndentState}{\State \hspace{1em}}
\newcommand{\DIndentState}{\State \hspace{2em}}

%% Temp
\newcommand{\disc}{\text{disc}}
\newcommand{\base}{\text{base}}

% Custom style for Python formatting
\lstset{
  language=Python,
  basicstyle=\ttfamily\small,
  keywordstyle=\color{blue},
  commentstyle=\color{gray},
  stringstyle=\color{purple},
  numbers=left,
  numberstyle=\tiny\color{gray},
  stepnumber=1,
  numbersep=10pt,
  backgroundcolor=\color{white},
  tabsize=4,
  showstringspaces=false,
  breaklines=true,
  breakatwhitespace=true
}

\title{EE588: Homework 3 Solutions}
\author{
  Naman Chhibbar \\
  USC ID: 7709-8548-02
}
\date{}

% Document body
\begin{document}
\maketitle

\section{Text exercises}

\begin{itemize}
  \item[4.15] \textit{Relaxation of Boolean LP.} In a \textit{Boolean Linear program}, the variable $x$ is constained to have components equal to zero or one:
  \begin{align*}
    \text{minimize} \quad & c^T x \\
    \text{subject to} \quad & Ax \preceq b \tag{4.67} \\
    & x_i \in \{0, 1\}, \quad i = 1, \ldots, n.
  \end{align*}
  In general, such problems are very difficult to solve, even though the feasible set is finite (containing at most $2^n$ points). \\
  In a general method called \textit{relaxation}, the constaint that $x_i$, has to be zero or one is replaced with the linear inequalities $0 \le x_i \le 1$:
  \begin{align*}
    \text{minimize} \quad & c^T x \\
    \text{subject to} \quad & Ax \preceq b \tag{4.68} \\
    & 0 \le x_i \le 1, \quad i = 1, \ldots, n.
  \end{align*}
  We refer to this problem as the \textit{LP relaxation} of the Boolean LP (4.67). The LP relaxation is far easier to solve than the original Boolean LP.
  \begin{itemize}
    \item[(a)] Show that the optimal value of the LP relaxation (4.68) is a lower bound on the optimal value of the Boolean LP (4.67). What can you say about the Boolean LP if the LP relaxation is infeasible?
    \item[(b)] It sometimes happens that the LP relaxation has a solution with $x_i \in \{0, 1\}$. What can you say in this case?
  \end{itemize}
  \textbf{Solution:}
  \begin{itemize}
    \item[(a)] Let $F_{\text{bool}}$ and $F_{\text{rel}}$ be the feasible sets of the Boolean LP and its LP relaxation, respectively. Now, since $\{0, 1\} \subset [0, 1]$, every point in $F_{\text{bool}}$ is also in $F_{\text{rel}}$, i.e., $F_{\text{bool}} \subseteq F_{\text{rel}}$. Since the optimal value of any minimization problem is the infimum of the objective function over the feasible set, we conclude that $p^*_{\text{bool}} \ge p^*_{\text{rel}}$, where $p^*_{\text{bool}}$ and $p^*_{\text{rel}}$ are the optimal values of the Boolean LP and its LP relaxation, respectively. \\
    If the LP relaxation is infeasible, it would mean that $F_{\text{rel}} = \emptyset$. Since $F_{\text{bool}} \subseteq F_{\text{rel}}$, we would also have $F_{\text{bool}} = \emptyset$, i.e., the Boolean LP is also infeasible.
    \item[(b)] Suppose $x^*_{\text{rel}}$ be an optimal solution to the LP relaxation such that $(x^*_{\text{rel}})_i \in \{0, 1\}$, i.e., $x^*_{\text{rel}} \in F_{\text{bool}}$. If $p^*_{\text{bool}}$ is the optimal value of the Boolean LP, then
    \[ p^*_{\text{bool}} \le c^T x^*_{\text{rel}} \]
    But, from part (a), we know that $p^*_{\text{bool}} \ge p^*_{\text{rel}} = c^T x^*_{\text{rel}}$. Hence, we get $p^*_{\text{bool}} = c^T x^*_{\text{rel}}$. Therefore, $x^*_{\text{rel}}$ is an optimal solution to the Boolean LP.
  \end{itemize}

  \item[4.17] \textit{Optimal activity levels.} We consider the selection of $n$ nonnegative activity levels, denoted $x_1, \ldots, x_n$. These activites consume $m$ resources, which are limited. Activity $j$ consumes $A_{ij}x_j$ of resource $i$, where $A_{ij}$ are given. The total resource consumption is additive, so the total of resource $i$ consumed is $c_i = \sum_{j=1}^n A_{ij} x_j$. (Ordinarily we have $A_{ij} \ge 0$, which means that activity $j$ consumes resource $i$. But we allow the possibility that $A_{ij} < 0$, which means that activity $j$ actually \textit{generates} resource $i$ as a by-product.) Each resource consumption is limited: we must have $c_i \le c_i^{\max}$, where $c_i^{\max}$ are given. Each activity generates revenuem which is a piecewise-linear concave function of the activity level:
  \[ r_j(x_j) = \begin{cases} p_j x_j & 0 \le x_j \le q_j \\ p_j q_j + p_j^\disc(x_j - q_j) & x_j \ge q_j \end{cases} \]
  Here $p_j > 0$ is the basic price, $q_j > 0$ is the quantity discount level, and $p_j^\disc$ is the quantity discount price, for (the product of) activity $j$. (We have $0 < p_j^\disc < p_j$.) The total revenue is the sum of the revenues from associated with each activity, i.e., $\sum_{j=1}^n r_j(x_j)$. The goal is to chose activity levels that maximize the total revenue while respecting the resource limits. Show how to formulate this problem as an LP. \\
  \textbf{Solution:} Since the revenue function $r_j$ is a piecewise-linear function, we can formulate an LP by splitting the activity levels $x_j$ into two parts:
  \begin{itemize}
    \item $x_j^\base$: The part of $x_j$ that generates revenue at the basic price $p_j$, hence $x_j^\base \le q_j$
    \item $x_j^\disc$: The part of $x_j$ that generates revenue the discounted price $p_j^\disc$, hence $x_j^\disc \ge 0$ only if $x_j^\base = q_j$, otherwise $x_j^\disc = 0$ (due to optimimzation reasons).
  \end{itemize}
  Therefore, $x_j = x_j^\base + x_j^\disc$. With this split, the LP formulation becomes much easier:
  \begin{align*}
    \text{maximize} \quad & \sum_{j=1}^n p_j x_j^\base + p_j^\disc x_j^\disc \\
    \text{subject to} \quad & \sum_{j=1}^n A_{ij} x_j \le c_j^{\max} \quad \forall i \in \{1, \ldots, m\} \\
    & 0 \le x_j^\base \le q_j \quad \forall j \in \{1, \ldots, n\} \\
    & 0 \le x_j^\disc \quad \forall j \in \{1, \ldots, n\}
  \end{align*}
  Again, this works since optmization would prefer increasing $x_j^\base$ till $x_j^\base = q_j$ since $p_j > p_j^\disc$. \\
  \textit{(The solution is supposed to end after formulating the LP, as requested in the question.)}

  \item[5.5] \textit{Dual of a general LP.} Find the dual function of the LP
  \begin{align*}
    \text{minimize} \quad & c^T x \\
    \text{subject to} \quad & Gx \preceq h \\
    & Ax = b.
  \end{align*}
  Give the dual problem, and make the implicit equality constraints explicit. \\
  \textbf{Solution:} Let $\lambda \preceq 0$ and $\nu$ be the dual vectors. The Lagrangian will be:
  \begin{align*}
    L(x, \lambda, \nu) &= c^T x + \lambda^T (Gx - h) + \nu^T (Ax - b) \\
    &= (c + G^T \lambda + A^T \nu)^T x - h^T \lambda - b^T \nu
  \end{align*}
  The Lagrange dual function will be:
  \[ g(\lambda, \nu) = \inf_x L(x, \lambda, \nu) = \inf_x \left( (c + G^T \lambda + A^T \nu)^T x - h^T \lambda - b^T \nu \right) \]
  The dual problem will be:
  \begin{align*}
    \text{maximize} \quad & g(\lambda, \nu) \\
    \text{subject to} \quad & \lambda \preceq 0
  \end{align*}
  After making the implicit dual constrain $c + G^T \lambda + A^T \nu = 0$ explicit, we get:
   \begin{align*}
    \text{maximize} \quad & -h^T \lambda - b^T \nu \\
    \text{subject to} \quad & \lambda \preceq 0 \\
    & c + G^T \lambda + A^T \nu = 0
  \end{align*}

  \item[5.13] \textit{Lagrangian relaxation of Boolean LP.} A \textit{Boolean linear program} is an optimization problem of the form
  \begin{align*}
    \text{minimize} \quad & c^T x \\
    \text{subject to} \quad & Ax \preceq b \\
    & x_i \in \{0, 1\}, \quad i = 1, \ldots, n,
  \end{align*}
  and is, in general, very difficult to solve. In exercise 4.15 we studies the LP relaxation of this problem,
  \begin{align*}
    \text{minimize} \quad & c^T x \\
    \text{subject to} \quad & Ax \preceq b \tag{5.107} \\
    & 0 \le x_i \le 1, \quad i = 1, \ldots, n,
  \end{align*}
  which is far easier to solve, and gives a lower bound on the optimal value of the Boolean LP. In this problem, we derive another lower bound for the Boolean LP, and work out the relation between the two lower bounds.
  \begin{itemize}
    \item[(a)] \textit{Lagrangian relaxation.} The Boolean LP can be reformulated as the problem
    \begin{align*}
      \text{minimize} \quad & c^T x \\
      \text{subject to} \quad & Ax \preceq b \\
      & x_i(1 - x_i) = 0, \quad i = 1, \ldots, n,
    \end{align*}
    which has quadratic equality constraints. Find the Langrange dual of this problem. The optimal value of the dual (which is convex) gives a lower bound on the optimal value of the Boolean LP. This method of finding a lower bound on the optimal value is called \textit{Lagrangian relaxation}.
    \item[(b)] Show that the lower bound obtained via Lagrangian relaxation, and via the LP relaxation (5.107), are the same. \textit{Hint.} Derive the dual of the LP relaxation (5.107). \\
    \textbf{Solution:}
    \begin{itemize}
      \item[(a)] With the dual variables $\lambda \succeq 0$ and $\nu$, we have the Lagrangian:
      \begin{align*}
        L(x, \lambda, \nu) &= c^T x + \lambda^T(Ax - b) + \sum_{i=1}^n \nu_i (x_i^2 - x_i) \\
        &= \sum_{i=1}^n \left( \nu_i x_i^2 + (c_i + (A^T \lambda)_i - \nu_i) x_i \right) - b^T \lambda
      \end{align*}
      The dual function is $g(\lambda, \nu) = \inf_x L(x, \lambda, \nu)$. We minimize over each $x_i$ independently:
      \begin{itemize}
        \item[1.] Case 1 $\nu_i < 0$: If a $\nu_i < 0$, there is a downwards parabola for an $i$ and $L \to -\infty$.
        \item[2.] Case 2 $\nu_i = 0$: $L$ is linear in $x_i$. The infimum is finite (and equals $0$) if and only if $c_i + (A^T \lambda)_i = 0$.
        \item[3.] Case 3 $\nu_i > 0$: $L$ is a strictly convex quadratic in $x_i$. Setting $\frac{\partial L}{\partial x_i} = 2\nu_i x_i + c_i + (A^T \lambda)_i - \nu_i = 0$ yields the minimizer:
        \[ x_i^* = \frac{\nu_i - c_i - (A^T \lambda)_i}{2\nu_i} \]
        Substituting $x_i^*$ back into $L$ gives the minimal value for component $i$:
        \[ -\frac{(c_i + (A^T \lambda)_i - \nu_i)^2}{4\nu_i} \]
        Thus, the Lagrange dual function is:
        \[ g(\lambda, \nu) = \begin{cases} -b^T \lambda - \displaystyle\frac{1}{4} \sum_{i=1}^n \frac{(c_i + (A^T \lambda)_i - \nu_i)^2}{\nu_i} & \text{if } \lambda \succeq 0, \, \nu \succeq 0 \\ -\infty & \text{otherwise} \end{cases} \]
        The optimal value of this convex optimization problem gives a lower bound on the Boolean LP.
      \end{itemize}
      \item[(b)] The LP relaxation is:
      \[ \begin{aligned} \text{minimize} \quad & c^T x \\ \text{subject to} \quad & Ax \preceq b \\ & -x \preceq 0 \\ & x \preceq 1 \end{aligned} \]
      Assign dual variables $u \succeq 0$ to $Ax \preceq b$, $v \succeq 0$ to $-x \preceq 0$, and $w \succeq 0$ to $x \preceq 1$. The Lagrangian for the LP relaxation will be:
      \begin{align*}
        L(x, u, v, w) &= c^T x + u^T(Ax - b) - v^T x + w^T(x - 1) \\
        &= (c + A^T u - v + w)^T x - b^T u - 1^T w
      \end{align*}
      Taking the infimum over $x$ yields the dual problem of the LP relaxation:
      \[ \begin{aligned} \text{maximize} \quad & -b^T u - 1^T w \\ \text{subject to} \quad & A^T u - v + w + c = 0 \\ & u \succeq 0, \quad v \succeq 0, \quad w \succeq 0 \end{aligned} \]
      From $A^T u - v + w + c = 0$, we have $w_i - v_i = -c_i - (A^T u)_i$.
      For a fixed $u$, we minimize $1^T w = \sum_{i=1}^n w_i$ over $v_i \ge 0, w_i \ge 0$:
      \begin{itemize}
        \item If $c_i + (A^T u)_i \ge 0$, setting $w_i = 0$ and $v_i = c_i + (A^T u)_i$ achieves the lower bound $w_i = 0$.
        \item If $c_i + (A^T u)_i < 0$, setting $w_i = -(c_i + (A^T u)_i)$ and $v_i = 0$ achieves $w_i = -(c_i + (A^T u)_i)$.
      \end{itemize}
      Thus, the optimal choice of $w_i$ is $\max(0, -(c_i + (A^T \lambda)_i))$.
      The maximized objective value of the LP dual w.r.t. $(v, w)$ is:
      \[ \phi(\lambda) = -b^T \lambda + \sum_{i=1}^n \min\left(0, \, c_i + (A^T \lambda)_i\right) \]
      Now consider the Lagrangian dual objective from part (a) for a fixed $\lambda$. We maximize over $\nu_i \ge 0$:
      \[ h(\nu_i) = -\frac{(a_i - \nu_i)^2}{4\nu_i}, \qquad \text{where } a_i = c_i + (A^T \lambda)_i \]
      Differentiating $h(\nu_i)$ with respect to $\nu_i \ge 0$ gives the optimal choice $\nu_i^* = \vert{}a_i\vert{} = \vert{}c_i + (A^T \lambda)_i\vert{}$:
      \begin{itemize}
        \item If $a_i \ge 0$, then $\nu_i^* = a_i \implies h(\nu_i^*) = 0$.
        \item If $a_i < 0$, then $\nu_i^* = -a_i \implies h(\nu_i^*) = a_i$.
      \end{itemize}
      In both cases, $\max_{\nu_i \ge 0} h(\nu_i) = \min(0, a_i) = \min(0, c_i + (A^T \lambda)_i)$. Hence, maximizing over the dual variables leads to the exact same reduced objective function $\phi(\lambda)$ in both cases.
    \end{itemize}
  \end{itemize}
\end{itemize}

\section{Additional exercises}

\begin{itemize}
  \item[1.] \textit{Reformulating constraints in CVX.} Each of the following CVX code fragments describes a convex constraint on the scalar variables $x$, $y$, and $z$, but violates the CVX rule set, and so is invalid. Briefly explain why each fragmet is invalid. Then, rewrite each one in an equivalent form that conforms to the CVX rule set. In your reformulations, you can use linear equality and inequality constraints, and inequalities constructed using CVX functions. You can also introudce additional variables, or use LMIs/ Be sure to explain (briefly) why your reformulation is equivalent to the original constraint, if it is not obvious. \\
  Check your reformulations by creating a small problem that includes these contraints, and solving it using CVX. Your test problems doesn't have to be feasible; it's enough to verify that CVX processes your constraints without error. \\
  \textit{Remark.} This looks like a problem about `how to use CVX software', or `tricks for using CVX'. But it really checks whether you understand the various composition rules, convex analysis, and constraint reformulation rules.
  \begin{itemize}
    \item[1.] $\mathtt{norm([x + 2*y, x - y]) == 0}$
    \item[2.] $\mathtt{square(square(x + y)) <= x - y}$
    \item[3.] $\mathtt{1/x + 1/y <= 1; x >= 0; y >= 0}$
    \item[4.] $\mathtt{norm([\max(x, 1), \max(y, 2)]) <= 3*x + y}$
    \item[5.] $\mathtt{x*y >= 1; x>= 0; y >= 0}$
    \item[6.] $\mathtt{(x + y)\caret{2}/sqrt(y) <= x - y + 5}$
    \item[7.] $\mathtt{x\caret{3} + y\caret{3} <= 1; x >= 0; y >= 0}$
    \item[8.] $\mathtt{x + z <= 1 + sqrt(x*y - z\caret{2}); x >= 0; y >= 0}$
  \end{itemize}
  \textbf{Solution:}
  \begin{itemize}
    \item[1.] This snippet violates the CVX rule set since an equality constraint (\texttt{==}) require both sides to be affine, whereas \texttt{norm} is not an affine. \\
    Since the constraint equals to zero, we can easily reformulate the snippet so that the vector elements are zero, which are affine:
    \begin{align*}
      \mathtt{x + 2*y == 0} \\
      \mathtt{x - y == 0}
    \end{align*}
    \item[2.] This violates DCP since in a composition $f(g(\ldots))$ with $f$ and $g$ convex, $f$ must be non-decreasing (under DCP rules). Since the outer function $f$ is also a $\mathtt{square}(u)$ function, CVX rejects it since it is only monotonic if $u \ge 0$. \\
    We can fix the same by introducing a new auxillary variable $t$ as follows: 
    \begin{align*}
      \mathtt{variable \quad t} \\
      \mathtt{square(x + y) <= t} \\
      \mathtt{square(t) <= x - y}
    \end{align*}
    This works since optimization would pull $t$ as close to $\mathtt{square}(x+y)$ as possible. It would do the same for $\mathtt{square}(t)$ by getting close to $x - y$. Essentially, we are getting $\mathtt{square(square(x + y)) <= x - y}$.
    \item[3.] This is disallowed since DCP rules disallow variables in dominator of ``$\mathtt{/}$''. Instead, we should use the built-in function $\mathtt{inv\_pos}$ to invert. The reformulation looks like:
    \[ \mathtt{inv\_pos(x) + inv\_pos(y) <= 1} \]
    with the other two constraints same.
    \item[4.] DCP rules require $\mathtt{norm}$ to have affine arguments. Here, $\max$ is a non-affine operation. We can replace it by two auxillary variables as:
    \begin{align*}
      \text{variables} \quad & u ~ v \\
      u >= x \\
      u >= 1 \\
      v >= y \\
      v >= 2
    \end{align*}
    Since $u \ge x$ and $u \ge 1$, we can confirm that optimization would pull $u$ to be $\max(x, 1)$, and the same with $v$.
    \item[5.] CVX requires the left-side of $\mathtt{>=}$ to be concave, whereas $\mathtt{x*y}$ is neither concave or convex. Valid reformulation would be:
    \[ \mathtt{inv\_pos(x) <= y} \]
    \item[6.] CVX does not allow dividing by a variable square root like $\mathtt{sqrt(y)}$. Instead, we introduce a new auxillary variable $t$ to be the square root of $y$ as follows:
    \begin{align*}
      \mathtt{variable \quad t} \\
      \mathtt{quad\_over\_lin(x + y, t) <= x - y + 5} \tag{1} \\
      \mathtt{square(t) <= y} \\
      t >= 0
    \end{align*}
    Since, $t$ is in the denominator in (1), and the fact that $t$ is allowed to increase to $\sqrt{y}$, makes up for the fact that optimization would pull $t$ to be $\sqrt{y}$. Hence, this reformulation is equivalent to the original problem. 
    \item[7.] CVX disallows the use of explicit power operator $\caret$ on decision variables since it is not convex on all of $\mathbb{R}$. Instead, we should use $\mathtt{pow\_pos}$ whenever decision variables are non-negative, as follows:
    \[ \mathtt{pow\_pos(x, 3) + pow\_pos(y, 3) <= 1} \]
    with the rest of the constraints as is.
    \item[8.] The snippet violates CVX rules since the right of $\mathtt{<=}$ contains a $\mathtt{sqrt}$ with an argument that is not concave ($x*y - z\caret2$), which is now allowed. We can instead use the built-in function:
    \[ \mathtt{det\_rootn(x, z; z, y)} = \sqrt{x*y - z^2} \]
    that achieves the same goal.
  \end{itemize}

  \item[2.] \textit{Optimal activity levels.} Solve the optimal activity level problem described in exercise 4.17 in \textit{Convex Optimization}, for the instance with problem data
  \[ A = \begin{bmatrix} 1 & 2 & 0 & 1 \\ 0 & 0 & 3 & 1 \\ 0 & 3 & 1 & 1 \\ 2 & 1 & 2 & 5 \\ 1 & 0 & 3 & 2 \end{bmatrix}, \quad
  c^{\max} = \begin{bmatrix} 100 \\ 100 \\ 100 \\ 100 \\ 100 \end{bmatrix}, \quad
  p = \begin{bmatrix} 3 \\ 2 \\ 7 \\ 6 \end{bmatrix}, \quad
  p^\disc = \begin{bmatrix} 2 \\ 1 \\ 4 \\ 2 \end{bmatrix}, \quad
  q = \begin{bmatrix} 4 \\ 10 \\ 5 \\ 10 \end{bmatrix} \]
  \textbf{Solution:} Here is the Python code used to calculate the output:
  \begin{lstlisting}
    import cvxpy as cp
    import numpy as np
    # Data
    A = np.array([
        [1, 2, 0, 1],
        [0, 0, 3, 1],
        [0, 3, 1, 1],
        [2, 1, 2, 5],
        [1, 0, 3, 2]
    ], dtype=float)
    c_max = np.array([100, 100, 100, 100, 100], dtype=float)
    p = np.array([3, 2, 7, 6], dtype=float)
    p_disc = np.array([2, 1, 4, 2], dtype=float)
    q = np.array([4, 10, 5, 10], dtype=float)
    # Decision variable: activity levels (x >= 0)
    x = cp.Variable(4, nonneg=True)
    # Elementwise concave revenue function: r_j = min(p_j*x_j, p_j*q_j + p_disc_j*(x_j - q_j))
    r = cp.minimum(cp.multiply(p, x), cp.multiply(p, q) + cp.multiply(p_disc, x - q))
    # Define and solve problem
    problem = cp.Problem(cp.Maximize(cp.sum(r)), [A @ x <= c_max])
    problem.solve()
    # Results
    x_opt = x.value
    r_opt = r.value
    avg_price = r_opt / x_opt
    print(f"Optimal activity levels (x): {np.round(x_opt, 4)}")
    print(f"Revenue per activity (r): {np.round(r_opt, 4)}")
    print(f"Average price per unit: {np.round(avg_price, 4)}")
    print(f"Total revenue: ${problem.value:.2f}")
  \end{lstlisting}
  And, here is the output and answer to the problem: \\
  \phantom{} \\
  Optimal activity levels (x): [ 4.  22.5 31.   1.5] \\
  Revenue per activity (r): [ 12.   32.5 139.    9. ] \\
  Average price per unit: [3.     1.4444 4.4839 6.    ] \\
  Total revenue: \$192.50
\end{itemize}

\end{document}
